\documentclass[12pt,a4paper]{article}

% ---------- PREAMBLE: packages and page setup ----------
\usepackage[margin=25mm]{geometry}
\usepackage[T1]{fontenc}
\usepackage{graphicx}
\usepackage{amsmath,amssymb}
\usepackage{array,tabularx,multirow}
\usepackage{enumitem}
\usepackage{multicol}
\usepackage{setspace}
\usepackage{xcolor}
\usepackage{makeidx}
\usepackage[colorlinks=true,linkcolor=blue,urlcolor=blue,citecolor=blue]{hyperref}

\makeindex
\setlength{\parskip}{5pt}
\setlength{\parindent}{0pt}

% ---------- TITLE INFORMATION ----------
\title{\textbf{Introduction to Artificial Intelligence}\\
       \large A LaTeX Article Demonstrating Syllabus Topics}
\author{Your Name\\Department of Computer Science}
\date{\today}

\begin{document}
\maketitle

\begin{abstract}
This sample article demonstrates document structure, formatting, lists,
tables, figures, mathematical equations, cross-references, citations,
hyperlinks, an index, and multi-column text.
\end{abstract}

\tableofcontents
\listoffigures
\listoftables
\newpage

\section{Introduction}\label{sec:intro}
Artificial Intelligence (AI) studies systems that perform tasks requiring
human-like abilities. This article is a practical demonstration of
LaTeX\index{LaTeX}, not a technical survey of AI.

\textbf{Bold text}, \textit{italic text}, \underline{underlined text},
\texttt{monospaced text}, and \emph{emphasised text} are examples of
formatting. A footnote can be added like this.\footnote{This is a sample
footnote.}

\subsection{Page and paragraph formatting}
The \texttt{geometry} package controls page margins. The
\texttt{setspace} package provides line spacing commands:
\begin{itemize}[label=\textbullet,leftmargin=2em]
  \item \texttt{\string\singlespacing}
  \item \texttt{\string\onehalfspacing}
  \item \texttt{\string\doublespacing}
\end{itemize}
\onehalfspacing
This paragraph uses one-and-a-half line spacing.
\singlespacing

\section{Lists}\label{sec:lists}
An unordered list:
\begin{itemize}
  \item Machine learning
  \item Computer vision
  \item Natural language processing
\end{itemize}

An ordered list:
\begin{enumerate}[label=\alph*)]
  \item Collect data
  \item Train a model
  \item Evaluate results
\end{enumerate}

A description list:
\begin{description}
  \item[Dataset] A collection of data.
  \item[Algorithm] A procedure for solving a problem.
\end{description}

\section{Tables}\label{sec:tables}
Table~\ref{tab:methods} demonstrates a caption, label, alignment,
horizontal rules, and a multirow cell.

\begin{table}[htbp]
\centering
\caption{Examples of AI tasks}\label{tab:methods}
\begin{tabular}{|l|l|c|}
\hline
\textbf{Area} & \textbf{Task} & \textbf{Example}\\
\hline
Machine learning & Classification & Spam detection\\
\hline
\multirow{2}{*}{Computer vision} & Image recognition & Face detection\\
\cline{2-3}
& Object detection & Road signs\\
\hline
\multicolumn{2}{|c|}{Total areas shown} & 2\\
\hline
\end{tabular}
\end{table}

A flexible-width table uses \texttt{tabularx}:
\begin{table}[htbp]
\centering
\caption{Flexible-width example}
\begin{tabularx}{0.9\linewidth}{|l|X|}
\hline
Topic & Explanation\\ \hline
AI & A field concerned with intelligent computer systems.\\ \hline
LaTeX & A document preparation system built on TeX.\\ \hline
\end{tabularx}
\end{table}

\section{Including Images}\label{sec:figures}
Figure~\ref{fig:workflow} demonstrates image insertion. Keep
\texttt{workflow.png} in the same folder as this source file.

\begin{figure}[htbp]
\centering
\includegraphics[width=0.8\linewidth]{workflow.png}
\caption{A simple input-process-output workflow}\label{fig:workflow}
\end{figure}

\section{Mathematical Expressions}\label{sec:math}
Inline mathematics looks like \(a^2+b^2=c^2\). Display mathematics:
\[
  E = mc^2
\]
A numbered equation:
\begin{equation}\label{eq:linear}
  y = mx + c
\end{equation}
An aligned derivation:
\begin{align}
  f(x) &= x^2 + 2x + 1\\
       &= (x+1)^2
\end{align}
A matrix:
\[
A=\begin{bmatrix}1&2\\3&4\end{bmatrix}
\]
A piecewise expression:
\[
f(x)=\begin{cases}
0,&x<0,\\
1,&x\geq 0.
\end{cases}
\]
See Equation~\eqref{eq:linear}.

\section{Columns of Text}\label{sec:columns}
\begin{multicols}{2}
LaTeX can also typeset prose in newspaper-style columns. This is
different from columns inside a table. The \texttt{multicol} package
provides the \texttt{multicols} environment. Text flows from one
column to the next automatically.
\end{multicols}

\section{Cross-References and Hyperlinks}\label{sec:refs}
See Section~\ref{sec:lists}, Table~\ref{tab:methods}, and
Figure~\ref{fig:workflow}. The page is \pageref{fig:workflow}.
Visit the \href{https://www.latex-project.org/}{LaTeX Project website}.
A visible URL can be printed with \url{https://www.overleaf.com/}.

\section{Citations and Bibliography}\label{sec:cites}
LaTeX is widely used for technical documents \cite{lamport1994}.
The bibliography database is stored separately in
\texttt{references.bib}.

\bibliographystyle{plain}
\bibliography{references}

\printindex
\end{document}
